\documentclass[10pt]{article}
\usepackage{amsmath,amssymb,geometry}
\geometry{margin=1in}
\title{A positive characteristic root disproves conjecture 00000008535}
\author{Local AI-assisted solution draft}
\date{3 October 2026}
\usepackage{url}
% Compact consistent title for a short mathematical note.
\makeatletter
\renewcommand{\maketitle}{\begin{center}
{\Large\bfseries\@title\par}\vspace{0.6em}
{\small\@author\par\@date}
\end{center}\vspace{0.5em}}
\makeatother
\setlength{\emergencystretch}{3em}
\begin{document}
\maketitle
\paragraph{Claim being disproved.}
The first clause asserts that every root of the characteristic polynomial of a supersolvable geometric lattice is real and negative. With the standard characteristic-polynomial convention, the two-element Boolean lattice already disproves negativity.
\paragraph{The lattice and its hypotheses.}
Let $L=\{0,1\}$ with $0<1$. Its meet and join are minimum and maximum; its rank function is $r(0)=0$, $r(1)=1$. It is atomistic: $1$ is its only atom, $0$ is the empty join of atoms, and $1$ is the join of itself. It is upper semimodular. Indeed, if $x$ covers $x\wedge y$, then necessarily $x=1$ and $y=0$, and $x\vee y=1$ covers $y=0$. Thus $L$ is a finite geometric lattice.

Both elements are modular: for $a\le b$, the identity $a\vee(m\wedge b)=(a\vee m)\wedge b$ holds for $m=0$ and $m=1$. Therefore the maximal chain $0<1$ consists entirely of modular elements, and $L$ is supersolvable.
\paragraph{Characteristic polynomial.}
The defining recurrence
\[
 \sum_{y\le x}\mu(0,y)=\begin{cases}1,&x=0,\\0,&x\ne0\end{cases}
\]
gives $\mu(0,0)=1$ and $\mu(0,1)=-1$. Hence
\[
 \chi_L(t)=\sum_{x\in L}\mu(0,x)t^{r(1)-r(x)}=t-1.
\]
It has the positive real root $t=1$, contrary to the asserted strict negativity. This suffices to disprove the whole conjecture, regardless of its additional root-modulus assertion.
\paragraph{Formal verification and semantic correspondence.}
The accompanying \texttt{Main.lean} is self-contained over Lean 4.19.0 \texttt{Std}. Its finite \texttt{LatticeData} contains an order, bottom, top, meet, join, and rank. It checks the partial-order and lattice universal properties, the grading of covers, atomisticity, and upper semimodularity. Supersolvability is expressed by a maximal chain all of whose elements satisfy the modular-element identity. These properties are proved for \texttt{booleanOne}.

The file checks the defining M\"obius recurrence for the actual two values and evaluates the characteristic polynomial by its defining sum. The theorem \path{positive_characteristic_root} proves that this sum vanishes at the integer $1$. The universal proposition \texttt{NegativeRootClaim} is the necessary integer-root specialization of the conjecture's first clause: every integer root of every supersolvable geometric lattice is negative. Every integer root is a real root under the canonical inclusion $\mathbb Z\subset\mathbb R$; therefore falsity of this specialization refutes the original real-root statement. No analytic assertion about real roots is needed.

The final theorem \texttt{conjecture8535\_false} negates \texttt{NegativeRootClaim} using all the verified hypotheses and the root $1$. Kernel-checked finite enumeration uses \texttt{decide}. The printed axiom lists contain only the standard axiom \texttt{propext}; there are no proof placeholders, external decision oracles, or additional axioms.
\paragraph{Reproduction.}
Run \texttt{lean Main.lean} using Lean 4.19.0.
\end{document}
